% from ega4-12-fr.tex
\oldpage[IV-3]{180}
  \item L'ensemble des $y\in Y$ tels que $\mathcal F_y$ ait la propriété
  $(S_k)$ et soit équidimensionnel en tout point de $X_y$.
  \item L'ensemble des $y\in Y$ tels que
  $\operatorname{coprof}(\mathcal F_y)\leq k$.
  \item L'ensemble des $y\in Y$ tels que $\mathcal F_y$ soit un
  $\mathcal O_{X_y}$-Module de Cohen-Macaulay.
  \item L'ensemble des $y\in Y$ tels que $\mathcal F_y$ soit géométriquement
  réduit.
  \item L'ensemble des $y\in Y$ tels que $\mathcal F_y$ soit géométriquement
  ponctuellement intègre en chaque point de $X_y$.
  \item L'ensemble des $y\in Y$ tels que $\mathcal F_y$ soit géométriquement
  intègre.
  \item L'ensemble des $y\in Y$ tels que $\mathcal F_y$ n'ait pas de cycle
  premier associé immergé et que la somme $l(y)$ des multiplicités totales
  (4.7.12) de $\mathcal F_y$ aux points maximaux de
  $\operatorname{Supp}(\mathcal F_y)$ soit $\leq k$.
\end{enumerate}
\end{theorem}

\begin{proof}
À l'exception de (ii), (iii), (iv), (x) et (xi), les propriétés $P(y)$
considérées sont de la forme suivante~: «~pour tout $x\in X_y$, $\mathcal F_y$
a la propriété $Q(x)$~», où $Q(x)$ est l'une des propriétés (i) à (viii) de
(12.1.1). Il résulte de (12.1.1) que l'ensemble $U$ des $x\in X$ tels que
$Q(x)$ soit vraie est ouvert, et l'ensemble $V$ des $y\in Y$ tels que $P(y)$
soit vraie n'est autre que l'ensemble $Y-f(X-U)$~; dans tous ces cas, le
théorème est donc déjà vrai en supposant seulement le morphisme $f$ \emph{fermé}.
Pour (iii), on applique (i) avec $\Phi$ réduit à un seul élément. Il reste à
examiner les cas (ii), (iv) et (xi) séparément, ((x) se déduisant aussitôt de
(xi) et de (4.7.14)) en appliquant toujours la méthode décrite dans (12.0.2).
\end{proof}

\begin{env}[12.2.1.1]
\label{IV.12.2.1.1-fr}
Cas (ii)~: Les étapes A) et B) de la démonstration procèdent comme dans le
début de (12.1.1)~; pour l'étape A), on utilise (8.9.1), (11.2.6) ainsi que
(4.2.7)~; pour l'étape B), on utilise (9.5.5) appliqué à
$\operatorname{Supp}(\mathcal F)$. Il reste à démontrer que si (en supposant
$Y$ noethérien) l'ensemble des dimensions des composantes irréductibles de
$\operatorname{Supp}(\mathcal F_y)$ contient $\Phi$, alors il en est de même
de l'ensemble des dimensions des composantes irréductibles de
$\operatorname{Supp}(\mathcal F_{y'})$, pour toute \emph{générisation} $y'$ de
$y$ dans $Y$. Soit $Y_1$ un spectre d'anneau de valuation discrète tel que, si
$y_1$ et $y'_1$ sont le point fermé et le point générique de $Y_1$, il y ait un
morphisme $h:Y_1\to Y$ tel que $h(y_1)=y$, $h(y'_1)=y'$ (\textbf{II}, 7.1.7).
Appliquant (4.2.7), on voit qu'on peut remplacer $Y$ par $Y_1$ et $X$ par
$X_1=X\times_YY_1$, autrement dit, se borner au cas où $Y$ est spectre
d'anneau de valuation discrète, $y$ son point fermé et $y'$ son point générique.

Utilisant (\textbf{Err}\textsubscript{\textbf{III}}, 30), on peut se borner au
cas où $\operatorname{Supp}(\mathcal F)=X$, de sorte que $f$ est quasi-plat
(2.3.3)~; les composantes irréductibles $Z_i$ de $X$ dominent alors $Y$
(2.3.4), autrement dit leurs points génériques appartiennent à $X_{y'}$, et
les $Z_i\cap X_{y'}$ sont les composantes irréductibles de $X_{y'}$. Mais toute
composante irréductible de $X_y$ est contenue dans un des $Z_i$, donc est une
composante irréductible de $(Z_i)_y$~; or, on sait (7.1.13) que les dimensions
de toutes les composantes irréductibles de $(Z_i)_y$ sont égales à
$\dim((Z_i)_{y'})$, ce qui achève de prouver (ii).
\end{env}

\begin{env}[12.2.1.2]
\label{IV.12.2.1.2-fr}
Cas (iv)~: Le même raisonnement qu'au début montre, en utilisant (12.1.1,
(iii)) que l'on peut déjà supposer (en remplaçant $Y$ par un ouvert de $Y$) que
pour tout $y\in Y$, $\mathcal F_y$ n'a pas de cycle premier associé immergé.
L'assertion du cas (iv) est alors une conséquence immédiate des assertions des
cas (i) et (ii).
\end{env}

\begin{env}[12.2.1.3]
\label{IV.12.2.1.3-fr}
